\chapter[Thème 1 $-$ Dérivation]{Thème 1 : Modèles définis par une fonction \\Dérivation}
\thispagestyle{fancy}

\section{Rappels sur la dérivation}

\subsection{Tangente à une courbe}

Soit une fonction $f$ définie sur un intervalle I et dérivable en un nombre réel $a$ appartenant à I.

$L$ est le nombre dérivé de $f$ en $a$.

A est un point d'abscisse $a$ appartenant à la courbe représentative $C_f$ de $f$.



    \begin{bdefi}
        \phantom{$\dfrac{1}{2}$}\pointilles[\textwidth-0.5cm]
        
        \phantom{$\dfrac{1}{2}$}\pointilles[\textwidth-0.5cm]
    \end{bdefi}
    
    \begin{bprop}
    	\phantom{$\dfrac{1}{2}$}\pointilles[\textwidth-0.5cm]
    	
    	\phantom{$\dfrac{1}{2}$}\pointilles[\textwidth-0.5cm]
    \end{bprop}
    
    \begin{bmethodet}{Déterminer l'équation d'une tangente}
    	\grasInMet{Enoncé :}
    	
    	On considère la fonction trinôme $f$ définie sur $\mathbb{R}$ par $f(x)=x^2+3 x-1$.
    	On veut déterminer une équation de la tangente à la courbe représentative de $f$ au point A de la courbe d'abscisse 2.
    	
    	\grasInMet{Réponse :}
    	
    	\phantom{$\dfrac{1}{2}$}\pointilles[\textwidth-0.5cm]
    	
    	\phantom{$\dfrac{1}{2}$}\pointilles[\textwidth-0.5cm]
    	
    	\phantom{$\dfrac{1}{2}$}\pointilles[\textwidth-0.5cm]
    	
    	\phantom{$\dfrac{1}{2}$}\pointilles[\textwidth-0.5cm]
    	
    	\phantom{$\dfrac{1}{2}$}\pointilles[\textwidth-0.5cm]
    	
    	\phantom{$\dfrac{1}{2}$}\pointilles[\textwidth-0.5cm]
    	
    	\phantom{$\dfrac{1}{2}$}\pointilles[\textwidth-0.5cm]
    \end{bmethodet}
    
    \subsection{formule de dérivation}
    
    \begin{minipage}[c]{0.48\linewidth}
    	\begin{center}
    		\renewcommand{\arraystretch}{2}
    		\begin{tabular}{|Xc|Xc|}
    			\hline Fonction & Dérivée \\
    			\hline $a, a \in \mathbb{R}$ & $\ldots$ \\
    			\hline $a x, a \in \mathbb{R}$ & $\ldots$ \\
    			\hline $x^2$ & $\ldots$ \\
    			\hline $x^n$, $n \geqslant 1$ entier & $\ldots$ \\
    			\hline $\dfrac{1}{x}$ & $\ldots$ \\
    			\hline $\dfrac{1}{x^n}$, $n \geqslant 1$ entier & $\ldots$ \\
    			\hline $\sqrt{x}$ & $\ldots$ \\
    			\hline $e^x$ & $\ldots$ \\
    			\hline $e^{k x}$, $k \in \mathbb{R}$ & $\ldots$ \\
    			\hline
    		\end{tabular}
    	\end{center}
    	
    \end{minipage} \hfill
    \begin{minipage}[c]{0.48\linewidth}
    	\begin{center}
    		\renewcommand{\arraystretch}{2.5}
    		\begin{tabular}{|c|c|}
    			\hline Fonction & Dérivée \\
    			\hline$u+v$ & $\ldots$ \\
    			\hline$k u$, $k \in \mathbb{R}$ & $\ldots$ \\
    			\hline$u v$ & $\ldots$ \\
    			\hline$\dfrac{1}{u}$ & $\ldots$ \\
    			\hline$\dfrac{u}{v}$ & $\ldots$ \\
    			\hline
    		\end{tabular}
    	\end{center}
    \end{minipage}
    
    \begin{bmethodet}{Déterminer la fonction dérivée}
    	\grasInMet{Enoncé :}
    	
    	Déterminer la fonction dérivée des fonctions suivantes:
    	\begin{multicols}{2}
    		\begin{enumMet}
    			\item $f(x)=\left(2 x^2-5 x\right)(3 x-2)$
    			\item $g(x)=\dfrac{6 x-5}{x^3-2 x^2-1}$
    		\end{enumMet}
    	\end{multicols}
    	\grasInMet{Réponse :}
    	
    	\begin{enumMet}
    		\item \phantom{$\dfrac{1}{2}$}\pointilles[\textwidth-0.5cm]
    		
    		\phantom{$\dfrac{1}{2}$}\pointilles[\textwidth-0.5cm]
    		
    		\phantom{$\dfrac{1}{2}$}\pointilles[\textwidth-0.5cm]
    		
    		\phantom{$\dfrac{1}{2}$}\pointilles[\textwidth-0.5cm]
    		
    		\phantom{$\dfrac{1}{2}$}\pointilles[\textwidth-0.5cm]
    		
    		\phantom{$\dfrac{1}{2}$}\pointilles[\textwidth-0.5cm]
    		
    		\phantom{$\dfrac{1}{2}$}\pointilles[\textwidth-0.5cm]
    		
    		\phantom{$\dfrac{1}{2}$}\pointilles[\textwidth-0.5cm]
    		\item \phantom{$\dfrac{1}{2}$}\pointilles[\textwidth-0.5cm]
    		
    		\phantom{$\dfrac{1}{2}$}\pointilles[\textwidth-0.5cm]
    		
    		\phantom{$\dfrac{1}{2}$}\pointilles[\textwidth-0.5cm]
    		
    		\phantom{$\dfrac{1}{2}$}\pointilles[\textwidth-0.5cm]
    		
    		\phantom{$\dfrac{1}{2}$}\pointilles[\textwidth-0.5cm]
    		
    		\phantom{$\dfrac{1}{2}$}\pointilles[\textwidth-0.5cm]
    		
    		\phantom{$\dfrac{1}{2}$}\pointilles[\textwidth-0.5cm]
    		
    		\phantom{$\dfrac{1}{2}$}\pointilles[\textwidth-0.5cm]
    	\end{enumMet}
    \end{bmethodet}
    
    \subsection{Application à l'étude des variations d'une fonction}
    
    \begin{bthm}
    	Soit une fonction $f$ définie et dérivable sur un intervalle $I$.
    	\begin{itemProp}
    		\item Si $f^{\prime}(x) \leqslant 0$, alors \phantom{$\dfrac{1}{2}$}\pointilles[\textwidth-5cm]
	    	\item Si $f^{\prime}(x) \geqslant 0$, alors \phantom{$\dfrac{1}{2}$}\pointilles[\textwidth-5cm]
    	\end{itemProp}
    \end{bthm}
    
    \begin{bmethodet}{Déterminer les variations d'une fonction}
    	\grasInMet{Enoncé :}
    	
    	Soit la fonction $f$ définie sur $\mathbb{R}$ par $f(x)=x^2-4 x$.
    	
    	Déterminer les variations de la fonction $f$.
    	
    	\grasInMet{Réponse :}
    	
    	\phantom{$\dfrac{1}{2}$}\pointilles[\textwidth-0.5cm]
    	
    	\phantom{$\dfrac{1}{2}$}\pointilles[\textwidth-0.5cm]
    	
    	\phantom{$\dfrac{1}{2}$}\pointilles[\textwidth-0.5cm]
    	
    	\phantom{$\dfrac{1}{2}$}\pointilles[\textwidth-0.5cm]
    	
    	\phantom{$\dfrac{1}{2}$}\pointilles[\textwidth-0.5cm]
    	
    	\phantom{$\dfrac{1}{2}$}\pointilles[\textwidth-0.5cm]
    	
    	\phantom{$\dfrac{1}{2}$}\pointilles[\textwidth-0.5cm]
    \end{bmethodet}
    
    \section{Dérivée d'une fonction composée}
    
    \begin{center}
    	\renewcommand{\arraystretch}{2}
    	\begin{tabular}{|Xc|Xc|}
    		\hline Fonction & Dérivée \\
    		\hline$f(a x+b)$ & $\ldots$ \\
    		\hline$u^2$ & $\ldots$ \\
    		\hline$e^u$ & $\ldots$ \\
    		\hline
    	\end{tabular}
    \end{center}
    
    \begin{bmethodet}{Déterminer la dérivée d'une fonction composée	}
    	\grasInMet{Enoncé :}
    	
    	Soit $f$ la fonction définie sur $\mathbb{R}$ par $f(x)=x e^{-\frac{x}{2}}$.
    	\begin{multicols}{2}
    		\begin{enumMet}
    			\item $f(x)=\left(2 x^2+3 x-3\right)^2$
    			\item $g(x)=2 e^{\frac{1}{x}}$
    		\end{enumMet}
    	\end{multicols}
    	
    	\grasInMet{Réponse :}
    	
    	\begin{enumMet}
    		\item \phantom{$\dfrac{1}{2}$}\pointilles[\textwidth-0.5cm]
    		
    		\phantom{$\dfrac{1}{2}$}\pointilles[\textwidth-0.5cm]
    		
    		\phantom{$\dfrac{1}{2}$}\pointilles[\textwidth-0.5cm]
    		
    		\phantom{$\dfrac{1}{2}$}\pointilles[\textwidth-0.5cm]
    		
    		\item \phantom{$\dfrac{1}{2}$}\pointilles[\textwidth-0.5cm]
    		
    		\phantom{$\dfrac{1}{2}$}\pointilles[\textwidth-0.5cm]
    		
    		\phantom{$\dfrac{1}{2}$}\pointilles[\textwidth-0.5cm]
    		
    		\phantom{$\dfrac{1}{2}$}\pointilles[\textwidth-0.5cm]
    	\end{enumMet}
    \end{bmethodet}
    
    \begin{bmethodet}{Étudier une fonction composée}
    	\grasInMet{Enoncé :}
    	
    	Soit $f$ la fonction définie sur $\mathbb{R}$ par $f(x)=x e^{-\frac{x}{2}}$.
   		\begin{enumMet}
   			\item Calculer la dérivée de la fonction $f$.
   			\item En déduire les variations de la fonction $f$.
   		\end{enumMet}
    	
    	\grasInMet{Réponse :}
    	
    	\begin{enumMet}
    		\item \phantom{$\dfrac{1}{2}$}\pointilles[\textwidth-0.5cm]
    		
    		\phantom{$\dfrac{1}{2}$}\pointilles[\textwidth-0.5cm]
    		
    		\phantom{$\dfrac{1}{2}$}\pointilles[\textwidth-0.5cm]
    		
    		\phantom{$\dfrac{1}{2}$}\pointilles[\textwidth-0.5cm]
    		
    		\item \phantom{$\dfrac{1}{2}$}\pointilles[\textwidth-0.5cm]
    		
    		\phantom{$\dfrac{1}{2}$}\pointilles[\textwidth-0.5cm]
    		
    		\phantom{$\dfrac{1}{2}$}\pointilles[\textwidth-0.5cm]
    		
    		\phantom{$\dfrac{1}{2}$}\pointilles[\textwidth-0.5cm]
    	\end{enumMet}
    \end{bmethodet}
    
    


